\documentclass[11pt]{article}
\usepackage{../../../tex/lecturestyle}

\title{Transposition Principle Introduction}
\author{}
\date{}

\begin{document}
\thispagestyle{fancy}

\lecturetitleband{Linear Algebra}{Transposition Principle Introduction}{6}

\section{Overview}
The transposition principle mechanically converts an algorithm for one
problem into an algorithm for a related \emph{transposed} problem ---
often a problem that looks quite different on its face. It is best
known for speeding up multipoint polynomial evaluation and for deriving
formal power series composition algorithms, but this lecture develops it
from scratch and applies it only to accessible, beginner-level examples,
to build intuition before those harder applications. The principle
(also called Tellegen's principle, after its origin in electrical
network theory) is usually presented in terms of signal-flow networks;
this lecture uses a simplified, purely algebraic presentation instead.
Assumes computing range sums via prefix sums; matrix familiarity helps
but is not required (the matrix-based proof and remarks can be skipped).

\section{The Transposed Problem}
\begin{definitionbox}[Problem 1(A)]
Given integers $A_{i,j}$ ($1\le i\le N$, $1\le j\le M$) and
$x_1,\ldots,x_M$, compute $y_1,\ldots,y_N$ where
$y_i = \sum_{j=1}^M A_{i,j}x_j$.
\end{definitionbox}
\begin{definitionbox}[Problem 1(B)]
Given the same $A_{i,j}$ and integers $Y_1,\ldots,Y_N$, compute
$X_1,\ldots,X_M$ where $X_j = \sum_{i=1}^N A_{i,j}Y_i$.
\end{definitionbox}
Problem 1(B) is the \emph{transposed problem} of Problem 1(A), and vice
versa.

\subsection{Restatement 1: Matrices}
With $x$ an $M\times1$ matrix, $y$ an $N\times1$ matrix, and $A$ an
$N\times M$ matrix, Problem 1(A) computes $y=Ax$; with $X$ an
$M\times1$ matrix and $Y$ an $N\times1$ matrix, Problem 1(B) computes
$X=A^TY$. So 1(A) multiplies a column vector by $A$, and 1(B)
multiplies a column vector by $A$'s transpose.

\subsection{Restatement 2: A Duality Identity}
The four quantities satisfy
\[
  \sum_{j=1}^M x_jX_j = \sum_{i=1}^N y_iY_i.
\]
So Problem 1(B) can be described as: given $Y$ (and knowing only that
$x,y$ are related as in Problem 1(A), without being given $x$ or $y$
themselves), find $X$ making this identity hold for every $x$. ($X_j$
and $Y_i$ are sometimes called \emph{dual variables} to $x_j$ and
$y_i$; this restatement is, in the language of dual maps and dual
bases, just a rephrasing of Restatement 1.)

\begin{examplebox}
For the map $y_1=2x_1+3x_2$, $y_2=4x_1+5x_2$, $y_3=6x_1+7x_2$, the
transposed problem is $X_1=2Y_1+4Y_2+6Y_3$, $X_2=3Y_1+5Y_2+7Y_3$.
Deriving this from Restatement 2: $X_1,X_2$ must satisfy
$x_1X_1+x_2X_2 = (2x_1+3x_2)Y_1+(4x_1+5x_2)Y_2+(6x_1+7x_2)Y_3$ for
every $x_1,x_2$; collecting the right side by $x_1$ and $x_2$ gives
exactly the stated $X_1,X_2$.
\end{examplebox}

\begin{notebox}
Not every ``compute $y$ from $x$'' problem has a meaningful transposed
problem this way --- the construction needs $y$ to depend
\emph{linearly} on $x$. Relations like $y_1=x_1^2$ or
$y_1=\min(x_1,x_2)$ do not fit this framework.
\end{notebox}

\section{The Transposition Principle}
\begin{theorembox}[Transposition Principle / Tellegen's Principle]
Suppose Problem 1(A) is solved by the algorithm:
\begin{enumerate}[leftmargin=1.7em]
  \item Initialize variables $z_1,\ldots,z_k$ to $0$.
  \item Load $x_1,\ldots,x_M$ into $z_{s_1},\ldots,z_{s_M}$.
  \item For $q=1,\ldots,Q$: add $c_q\times z_{b_q}$ to $z_{a_q}$.
  \item Read $y_1,\ldots,y_N$ out of $z_{t_1},\ldots,z_{t_N}$.
\end{enumerate}
Then Problem 1(B) is solved by:
\begin{enumerate}[leftmargin=1.7em]
  \item Initialize variables $Z_1,\ldots,Z_k$ to $0$.
  \item Load $Y_1,\ldots,Y_N$ into $Z_{t_1},\ldots,Z_{t_N}$.
  \item For $q=Q,Q-1,\ldots,1$: add $c_q\times Z_{a_q}$ to $Z_{b_q}$.
  \item Read $X_1,\ldots,X_M$ out of $Z_{s_1},\ldots,Z_{s_M}$.
\end{enumerate}
\end{theorembox}
In short: reverse the order of the elementary steps, swap each step's
two roles ($a_q\leftrightarrow b_q$), and swap the roles of the
load/store positions.

\subsection{Proof via Matrices}
Each part of the Problem 1(A) algorithm is a matrix: loading $x$ into
$z$ is $z=Sx$ for the $k\times M$ matrix $S$ with a $1$ in position
$(s_j,j)$ and $0$ elsewhere (summing contributions when some $s_j$
coincide, exactly as ordinary matrix multiplication would); step $q$
replaces $z$ with $A_qz$ for the $k\times k$ matrix $A_q$ that is the
identity except for a $c_q$ in position $(a_q,b_q)$; and reading $y$
from $z$ is $y=Tz$ for the $N\times k$ matrix $T$ with a $1$ in
position $(i,t_i)$. So the whole algorithm computes
$y = T\cdot A_QA_{Q-1}\cdots A_1\cdot S\cdot x$, i.e.\
$A = T\cdot A_QA_{Q-1}\cdots A_1\cdot S$. Using $(BC)^T=C^TB^T$
repeatedly,
\[
  A^T = S^T\cdot A_1^TA_2^T\cdots A_{Q-1}^TA_Q^T\cdot T^T,
\]
so $X=A^TY$ is computed by applying, in order, $T^T$ (reading $Y$ into
$Z$ at the positions $t_i$ --- summing when they coincide, since $T^T$
is exactly the loading-type matrix for this step), then
$A_1^T,\ldots,A_Q^T$ (each just $A_q$ with $a_q,b_q$ swapped), then
$S^T$ (reading $X$ out of $Z$ at positions $s_j$). This is precisely
the stated algorithm for Problem 1(B).

\subsection{Proof via the Duality Identity}
Comparing the two algorithms, Problem 1(B)'s steps retrace Problem
1(A)'s in reverse. By Restatement 2, it suffices to show
$\sum_jx_jX_j=\sum_iy_iY_i$ where $X$ is the output of the Problem
1(B) algorithm above. This follows from three easy observations:
$\sum_jx_jX_j=\sum_kz_kZ_k$ at the moment $x$ is loaded into $z$ and
$Y$ into $Z$ (both sides are $0$ except at the loaded positions, where
they agree termwise); $\sum_kz_kZ_k$ is unchanged by each pair of
steps $q$ (forward step $q$ on $z$, matched with backward step $q$ on
$Z$), since
$(z_a+cz_b)Z_a + z_bZ_b = z_aZ_a + z_b(Z_b+cZ_a)$; and
$\sum_kz_kZ_k=\sum_iy_iY_i$ at the moment $y$ is read from $z$ and $Y$
was loaded into $Z$, by the same termwise argument as the first step.

\section{Worked Example: Range-Sum Queries}
\begin{definitionbox}[Problem 3(A)]
Given a length-$N$ sequence $x=(x_0,\ldots,x_{N-1})$, answer $Q$
queries; query $i$ gives $L_i,R_i$ ($0\le L_i\le R_i\le N-1$) and asks
for $y_i=\sum_{j=L_i}^{R_i}x_j$.
\end{definitionbox}
\begin{definitionbox}[Problem 3(B), the transpose]
Given $Q$ triples $(L_i,R_i,Y_i)$, compute
$X_0,\ldots,X_{N-1}$ where $X_j$ is the sum of $Y_i$ over all $i$ with
$L_i\le j\le R_i$.
\end{definitionbox}
Problem 3(A) answers range-sum queries; its transpose, Problem 3(B),
answers range-add queries. The classic prefix-sum solution to 3(A) can
be mechanically transposed into a solution to 3(B).

\subsection{Solving 3(A)}
Initialize $s_0,\ldots,s_N,y_1,\ldots,y_Q$ to $0$; load $x_j$ into
$s_{j+1}$; for $j=0,\ldots,N-1$, add $s_j$ into $s_{j+1}$ (building
prefix sums); for $i=1,\ldots,Q$, add $s_{R_i+1}$ into $y_i$ and add
$-s_{L_i}$ into $y_i$; output $y_1,\ldots,y_Q$.

\subsection{Solving 3(B) via the Transposition Principle}
Applying Theorem 2 mechanically: initialize $S_0,\ldots,S_N,Y_1,\ldots,
Y_Q$ to $0$; load the given $Y_i$; for $i=1,\ldots,Q$ (any order, since
these steps don't depend on each other), add $-Y_i$ into $S_{L_i}$ and
add $Y_i$ into $S_{R_i+1}$; for $j=N-1,\ldots,1,0$ (reverse of 3(A)'s
forward loop), add $S_{j+1}$ into $S_j$; output $X_j=S_{j+1}$ for
$j=0,\ldots,N-1$.

\begin{notebox}
The most familiar solution to Problem 3(B) is probably the \emph{imos
method}: load $Y_i$ into $S_{L_i}$ (not $-Y_i$) and $-Y_i$ into
$S_{R_i+1}$ (not $+Y_i$), run the prefix-sum loop forward
($j=0,\ldots,N-1$, adding $S_j$ into $S_{j+1}$), and read
$X_j=S_j$ (not $S_{j+1}$). This differs from the mechanically-derived
solution above in loop direction and sign convention, but both compute
the same $X$. Transposing the imos method's own algorithm gives yet
another (loop-order-shuffled) solution to Problem 3(A) --- worth
trying by hand.
\end{notebox}

% FIGURE: tp_networks
\begin{figure}[H]
\centering
\begin{tikzpicture}[x=1cm,y=1cm]
\node[lvert,minimum size=7mm] (Fs0) at (0.00,0.00) {$s_0$};
\node[lvert,minimum size=7mm] (Fs1) at (2.20,0.00) {$s_1$};
\node[lvert,minimum size=7mm] (Fs2) at (4.40,0.00) {$s_2$};
\node[lvert,minimum size=7mm] (Fs3) at (6.60,0.00) {$s_3$};
\draw[larrow] (Fs0) -- (Fs1);
\draw[larrow] (Fs1) -- (Fs2);
\draw[larrow] (Fs2) -- (Fs3);
\node[lnode,minimum width=0.9cm] (Fx0) at (2.20,1.35) {$x_0$};
\draw[larrow,lectureaccent] (Fx0) -- (Fs1);
\node[lnode,minimum width=0.9cm] (Fx1) at (4.40,1.35) {$x_1$};
\draw[larrow,lectureaccent] (Fx1) -- (Fs2);
\node[lnode,minimum width=0.9cm] (Fx2) at (6.60,1.35) {$x_2$};
\draw[larrow,lectureaccent] (Fx2) -- (Fs3);
\node[lnode,minimum width=0.9cm,lhi2] (Fy) at (8.90,0.00) {$y$};
\draw[larrow,lecturewarn] (Fs3) -- node[above,llabel] {$+1$} (Fy);
\draw[larrow,lecturewarn,bend left=28] (Fs1) to node[above,llabel] {$-1$} (Fy);
\node[lvert,minimum size=7mm] (Ts0) at (0.00,-3.60) {$S_0$};
\node[lvert,minimum size=7mm] (Ts1) at (2.20,-3.60) {$S_1$};
\node[lvert,minimum size=7mm] (Ts2) at (4.40,-3.60) {$S_2$};
\node[lvert,minimum size=7mm] (Ts3) at (6.60,-3.60) {$S_3$};
\draw[larrow] (Ts1) -- (Ts0);
\draw[larrow] (Ts2) -- (Ts1);
\draw[larrow] (Ts3) -- (Ts2);
\node[lnode,minimum width=0.9cm] (Tx0) at (2.20,-2.25) {$X_0$};
\draw[larrow,lectureaccent] (Ts1) -- (Tx0);
\node[lnode,minimum width=0.9cm] (Tx1) at (4.40,-2.25) {$X_1$};
\draw[larrow,lectureaccent] (Ts2) -- (Tx1);
\node[lnode,minimum width=0.9cm] (Tx2) at (6.60,-2.25) {$X_2$};
\draw[larrow,lectureaccent] (Ts3) -- (Tx2);
\node[lnode,minimum width=0.9cm,lhi2] (Ty) at (8.90,-3.60) {$Y$};
\draw[larrow,lecturewarn] (Ty) -- node[above,llabel] {$+1$} (Ts3);
\draw[larrow,lecturewarn,bend left=28] (Ty) to node[above,llabel] {$-1$} (Ts1);
\node[llabel,anchor=east,align=right] at (-0.7,0.6) {(a) prefix sums:\\ $y=s_3-s_1=x_1+x_2$};
\node[llabel,anchor=east,align=right] at (-0.7,-3.0) {(b) its transpose:\\ range add of $Y$\\ on $[1,3)$};
\end{tikzpicture}
\caption{The range-sum network (a) for the single query $[1,3)$ over $n=3$ values, and its transpose (b) obtained by reversing every arrow: prefix-sum accumulation $s_j\to s_{j+1}$ becomes the suffix accumulation $S_{j+1}\to S_j$, and the query reads $+s_3,-s_1$ becomes the writes $+Y\to S_3,\ -Y\to S_1$. Pushing $Y=1$ through (b) gives $S_3=1,S_2=1,S_1=0,S_0=0$, hence $X=(S_1,S_2,S_3)=(0,1,1)$: the value $Y$ is added exactly to positions $1$ and $2$.}\label{fig:tp_networks}
\end{figure}

\section{Practice Problems}
\begin{examplebox}
\textbf{Subset zeta transform.} Given $x=(x_0,\ldots,x_{2^N-1})$ with
$y_i=\sum_{S\subseteq i}x_S$, the transposed problem is the
\emph{superset} zeta transform, $X_i=\sum_{S\supseteq i}Y_S$: the fast
subset zeta transform's elementary steps (in the Zeta \& M\"obius
Transform lecture) mechanically transpose into the fast superset zeta
transform.

\textbf{Divisor zeta transform.} Given $x=(x_1,\ldots,x_N)$ with
$y_i=\sum_{j\mid i}x_j$, the transposed problem is the \emph{multiple}
zeta transform, $X_i=\sum_{i\mid j}Y_j$: the fast divisor-relation zeta
transform's steps mechanically transpose into the fast
multiple-relation zeta transform.

\textbf{Point update, range sum.} Given a sequence supporting point
updates ($x_i\mathrel{+}=v$) and range-sum queries
($\sum_{j=L}^{R}x_j$), transposing a Fenwick-tree or segment-tree
solution mechanically produces a solution to the transposed problem ---
range-update, point-query --- i.e.\ the well-known \emph{dual}
Fenwick tree / dual segment tree technique.
\end{examplebox}

\section{Reference Implementations}
The examples below implement a generic interpreter for the small
``elementary-step'' programs described in Theorem 2, running both the
forward algorithm and its mechanically-derived transpose. They validate
the transposition principle itself by comparing the transpose's output
against $A^TY$ computed directly from the matrix $A$ that the forward
algorithm implements (recovered by running it on every standard basis
vector), over many random programs; reproduce the lecture's own small
worked example exactly; encode the prefix-sum solution to range-sum
queries and confirm its mechanically-derived transpose agrees with both
hand-derived solutions to range-add queries from the lecture (the
direct derivation and the imos method) and with a brute-force
definition; and encode the fast subset zeta transform, confirming its
mechanically-derived transpose matches a brute-force superset zeta
transform.

\codenote
\subsection{Rust}
\corelisting{Rust}{example.rs}

\subsection{Python}
\corelisting{Python}{example.py}

\section{Worked Example}
\begin{examplebox}
The transposition principle itself (transpose output equalling $A^TY$)
is confirmed exactly across 300 random small programs. The lecture's
own $2\times3$ worked example is reproduced exactly, both forward and
transposed, across 50 random input trials. For range-sum/range-add
queries, the mechanically-derived transpose, the lecture's own
hand-derived transpose, the imos method, and a brute-force definition
all agree exactly across 100 random trials (sequence lengths up to
$20$, up to $14$ queries). For the subset zeta transform, its
mechanically-derived transpose matches a brute-force superset zeta
transform exactly across 30 random trials (up to $6$ bits).
\end{examplebox}

\section{Further Reading}
\begin{itemize}[leftmargin=1.5em]
  \item AtCoder Algorithm Lectures index (source of this curriculum):
        \url{https://info.atcoder.jp/entry/algorithm_lectures/index}
  \item AtCoder ARC185 E, ``Adjacent GCD'' (an application of
        transposing a linear-combination computation to get per-term
        contribution coefficients):
        \url{https://atcoder.jp/contests/arc185/tasks/arc185_e}
  \item Bostan, Lecerf, Schost, ``Tellegen's Principle into Practice,''
        ISSAC 2003:
        \url{https://doi.org/10.1145/860854.860870}
\end{itemize}

\end{document}
